\paragraph{A finite-\(N\) matrix toy for the class relations.}
On the model space of the class above, take the base to be the unit two-sphere with coframe \(\theta^1=\mathrm d\theta\), \(\theta^2=\sin\theta\,\mathrm d\varphi\), orthonormal frame \(e_1=\partial_\theta\), \(e_2=(\sin\theta)^{-1}\partial_\varphi\), and the circle group as structure group, with the Levi-Civita connection \(\omega^{12}=-\cos\theta\,\mathrm d\varphi\). The class data of this connection are \(\mathsf T_{12}{}^2=\cot\theta\), from \([e_1,e_2]=-\cot\theta\,e_2\), and \(\mathsf R_{12}{}^{12}=-1\) with \(\mathcal S_{12}\) the fundamental vertical field \(\xi\) (the curvature two-form \(\mathrm d\omega^{12}\) evaluates to \(+1\) on \((e_1,e_2)\)), with \(\mathsf F=0\). The horizontal lifts \(\mathbb D_a=e_a-\omega_a\xi\) satisfy exactly \([\mathbb D_1,\mathbb D_2]=-\cot\theta\,\mathbb D_2-\xi\) and \([\lambda\xi,\mathbb D_a]=-(\mathbb D_a\lambda)\xi\) for every real \(\lambda\) on the base: the vertical part of \([\mathbb D_1,\mathbb D_2]\) alone is \(-\csc^2\theta\,\xi\neq0\), so the commutator of two first-order operators is not itself first-order when the curvature is nonzero, and the zeroth-order slot \(\tfrac12\mathsf R_{ab}{}^{cd}\mathcal S_{cd}\) of the class relation is not optional. At finite \(N\) the derivation part of the same structure is carried by matrices: with spin-\(j\) generators \(J_i\), \(N=2j+1\), radius-one coordinates \(X_i=J_i/\sqrt{j(j+1)}\), so that \([X_i,X_j]=\mathrm i\hbar_{\mathrm{eff}}\varepsilon_{ijk}X_k\) with \(\hbar_{\mathrm{eff}}=1/\sqrt{j(j+1)}\), the inner derivations \(\mathcal L_a=\hbar_{\mathrm{eff}}^{-1}\operatorname{ad}_{X_a}\) obey the Leibniz rule and close the rotation algebra exactly, \([\mathcal L_i,\mathcal L_j]=\mathrm i\varepsilon_{ijk}\mathcal L_k\) and \(\mathcal L_i(X_j)=\mathrm i\varepsilon_{ijk}X_k\), for every \(N\); moving coefficient functions across the noncommutative product errs from the Poisson bracket at order \(O(\hbar_{\mathrm{eff}})\), and the frame redundancy \(\sum_i x_i\rho_i=0\) of the sphere frame is violated at the same order. The ghost operator \(\sum_i\operatorname{ad}_{X_i}^2/\hbar_{\mathrm{eff}}^2\), with the Euclidean contraction, has, on \(\operatorname{End}(\mathbb C^N)\), the spectrum \(\{l(l+1):0\le l\le N-1\}\) with multiplicity \(2l+1\) exactly, so it reproduces the Laplace--Beltrami spectrum of the sphere up to the removal of the modes \(l>N-1\). This is the finite-size regularization of curved-geometry matrix configurations of \cite{HattoriSekoTsuchiya2026}, read on the class of this section. The toy is compact in base and fibre, Abelian in fibre, torsion-free and two-dimensional: it retains neither \(\mathrm{SO}(3,1)\) nor the noncompact fibre nor the \((7,3)\) kinetic signature nor the selection of the class by the matrix equation, and it discharges neither Assumption~\ref{ass:frame-lock} nor Assumption~\ref{ass:nondeg} nor Assumption~\ref{ass:local-lorentz}. Against the criterion of Chapter~\ref{ch:agenda} this is a conditional sub-lemma on a compact toy, and the status recorded for Problem~D above is unchanged.

\paragraph{The locked frame as a gauge-invariant composite.}
Elitzur's theorem \cite{Elitzur1975} rules out a nonzero expectation value of a gauge-variant local field for a compact structure group, and Fradkin and Shenker \cite{FradkinShenker1979} showed that the Higgs and confinement regimes of such theories are analytically connected without any local order parameter. Both statements are established for compact gauge groups on the lattice; their extension to the noncompact frame group is not established. The expectation value of Equation~\eqref{eq:VEV} is gauge-covariant under Equation~\eqref{eq:Phi-transform}, so on the gauge-invariant reading adopted here the locking of Assumption~\ref{ass:frame-lock} is to be represented not by the nonzero expectation value itself but by composites that are bilocal in the base, in the spirit of the bilocal fields of \cite{HoKawaiSteinacker2025}, with parallel transporters built from the connection of the class above relating the two points, and the composite metric of Equation~\eqref{eq:g-from-e} is invariant under tangent-space rotations. On this reading Assumption~\ref{ass:nondeg} becomes a non-degeneracy statement about the two-point composite rather than a consequence of a broken symmetry, and the locking is recorded as a complementarity of regimes rather than a local spontaneous breaking. The construction presupposes the connection content that the closure of Problem~D1 in Chapter~\ref{ch:agenda} requires to be specified, and is recorded as a mechanism whose selection by the matrix dynamics is open, with no change to the status of Problem~D above.
